\documentclass[11pt, a4paper]{article}

\usepackage[T1]{fontenc}
\usepackage[utf8]{inputenc}
\usepackage[a4paper, top=2.5cm, bottom=2.5cm, left=2.5cm, right=2.5cm]{geometry}
\usepackage{booktabs}
\usepackage{enumitem}
\usepackage[hidelinks]{hyperref}
\usepackage{parskip}
\usepackage{amsmath}
\usepackage{listings}
\usepackage{xcolor}

\definecolor{codigobg}{RGB}{245,245,245}
\definecolor{keyword}{RGB}{0,100,150}
\lstset{language=R, basicstyle=\ttfamily\small, backgroundcolor=\color{codigobg},
  keywordstyle=\color{keyword}\bfseries, frame=single, breaklines=true, showstringspaces=false}

\setlength{\parindent}{0pt}

\pagestyle{plain}

\begin{document}

\begin{center}
{\large\bfseries PONTIFICIA UNIVERSIDAD JAVERIANA}\\[3pt]
Facultad de Ciencias Econ\'omicas y Administrativas\\[2pt]
Estad\'istica 2026-II --- Prof.\ Jaime Polanco Jim\'enez\\[6pt]
{\LARGE\bfseries Problem Set 2}\\[3pt]
{\large Correlaci\'on y Paradoja de Simpson}\\[6pt]
Asignado: Agosto 14 (Sesi\'on 2) \quad$|$\quad Entrega: Agosto 21 (Sesi\'on 3)
\end{center}

\vspace{6pt}
\rule{\linewidth}{0.5pt}
\vspace{2pt}
{\centering\small Repositorio: \url{https://github.com/polanco-jaime/Curso-Estadistica-2026-3}\par}
\vspace{6pt}

\textbf{Instrucciones.} Este Problem Set es individual y se resuelve en R.
La entrega es digital al inicio de la clase indicada.
El c\'odigo R debe incluirse.

\section*{Datos}

Los datos se encuentran en el repositorio del curso en formato Parquet:

\begin{lstlisting}
library(arrow)
saber_pro <- read_parquet("datos/saber_pro/saber_pro.parquet")
\end{lstlisting}

\textbf{Diccionario de variables utilizadas:}

\vspace{4pt}
\begin{tabular}{l l l}
\toprule
\textbf{Variable} & \textbf{Descripci\'on} & \textbf{Tipo} \\
\midrule
\texttt{MOD\_INGLES\_PUNT}           & Puntaje m\'odulo de ingl\'es        & Num\'erico \\
\texttt{MOD\_RAZONA\_CUANTITAT\_PUNT} & Puntaje razonamiento cuantitativo  & Num\'erico \\
\texttt{MOD\_LECTURA\_CRITICA\_PUNT}  & Puntaje lectura cr\'itica           & Num\'erico \\
\texttt{MOD\_COMPETEN\_CIUDADA\_PUNT} & Puntaje competencias ciudadanas    & Num\'erico \\
\texttt{MOD\_COMUNICA\_ESCRITA\_PUNT} & Puntaje comunicaci\'on escrita      & Num\'erico \\
\texttt{MOD\_INGLES\_DESEM}           & Nivel MCER (A$-$, A1, A2, B1, B+) & Categ\'orico \\
\texttt{ESTU\_GENERO}                & G\'enero del estudiante              & Categ\'orico \\
\texttt{INST\_ORIGEN}                & Tipo de IES: Oficial / Privada     & Categ\'orico \\
\texttt{FAMI\_ESTRATOVIVIENDA}       & Estrato socioecon\'omico (1--6)     & Ordinal \\
\bottomrule
\end{tabular}

\section*{Enunciado}

\begin{enumerate}[label=\alph*)]
  \item \textbf{Matriz de correlaciones:} usando los datos de Saber Pro,
        calcule la matriz de correlaciones de
        Pearson y de Spearman para los 5 m\'odulos:
        \begin{itemize}
          \item \texttt{MOD\_INGLES\_PUNT}
          \item \texttt{MOD\_RAZONA\_CUANTITAT\_PUNT}
          \item \texttt{MOD\_LECTURA\_CRITICA\_PUNT}
          \item \texttt{MOD\_COMPETEN\_CIUDADA\_PUNT}
          \item \texttt{MOD\_COMUNI\_ESCRITA\_PUNT}
        \end{itemize}
        Visualice ambas matrices con \texttt{corrplot} o un heatmap.
        ¿Qu\'e pares de m\'odulos tienen mayor y menor correlaci\'on?
        ¿Hay diferencias importantes entre Pearson y Spearman? ¿Por qu\'e?

  \item \textbf{Tabla de contingencia:} construya una tabla de contingencia
        cruzando g\'enero (\texttt{ESTU\_GENERO}) con nivel de ingl\'es seg\'un el MCER
        (A$-$, A1, A2, B1, B+). Calcule las frecuencias absolutas y relativas.
        ¿Hay alg\'un patr\'on visible?

  \item \textbf{Paradoja de Simpson:} demuestre la Paradoja de Simpson con datos
        ICFES. Busque una relaci\'on entre dos variables que cambie de direcci\'on
        (o desaparezca) al controlar por una tercera variable.

        \textit{Sugerencia:} compare la correlaci\'on entre \texttt{MOD\_INGLES\_PUNT}
        y otra variable a nivel agregado vs.\ desagregado por tipo de IES
        (p\'ublica/privada) o por estrato.

        Presente gr\'aficos que ilustren claramente c\'omo la relaci\'on se invierte
        o modifica al condicionar.

  \item \textbf{Conclusiones:} en un p\'arrafo, explique por qu\'e la Paradoja de
        Simpson es relevante para el an\'alisis de pol\'itica educativa. ¿Qu\'e
        riesgos implica no controlar por variables confusoras?
\end{enumerate}

\vspace{10pt}
\rule{\linewidth}{0.4pt}
\begin{center}
{\small Cada entrega completa suma $+0.0625$ a la nota final.}
\end{center}

\end{document}
